\documentclass[11pt,letterpaper]{article}
\usepackage{technical-user-manual}
\setDocTitle{Continuous Testing Architecture and Flaky-Test Triage}
\setDocSubtitle{Engineering technical reference}
\setDocVersion{1.0}
\setDocStatus{Technical education / proposed reference}
\setDocOwner{CI/CD Engineering}
\setDocDate{30 October 2026}
\setDocSummary{Supplementary engineering practice}
\begin{document}
\makedoctitle
\makedocsummary
\tableofcontents
\clearpage
\section{Purpose, scope and authority}
This is an original modern engineering reference. The 2007 Duvall, Matyas and Glover book informs foundational CI intent but does not prescribe GitHub-specific implementations or 2026 supply-chain standards.

\section{Portfolio}
Unit tests isolate business rules; component tests exercise module interfaces; contract tests detect consumer/provider drift; integration tests exercise databases and external boundaries; end-to-end tests validate selected user journeys and are slower and less diagnostic.

Verification: collect the configuration or runtime evidence supporting this control, identify an accountable owner, and record the exception path where the condition cannot presently be met.

\section{Isolation}
Each test owns temporary files, database state, time and randomness. Network access, ordering assumptions and global shared fixtures must be explicit. Run in shuffled order where practical; collect seed and container image digest.

Verification: collect the configuration or runtime evidence supporting this control, identify an accountable owner, and record the exception path where the condition cannot presently be met.

\section{Flakiness detection}
Re-run failed tests under same commit, environment and recorded seed to distinguish nondeterminism from changes. Quarantine requires a linked issue, owner, maximum duration and a replacement risk control; do not mark an unstable critical gate as optional indefinitely.

Verification: collect the configuration or runtime evidence supporting this control, identify an accountable owner, and record the exception path where the condition cannot presently be met.

\section{Coverage interpretation}
Coverage is a discovery aid for untested code; it does not demonstrate assertion strength or absence of defects. Analyze mutation score or intentionally fault-injected regressions where appropriate. Treat service-level SLO and load-test evidence as separate signals.

Verification: collect the configuration or runtime evidence supporting this control, identify an accountable owner, and record the exception path where the condition cannot presently be met.

\section{Reporting}
Publish stable JUnit XML, coverage and relevant logs with retention and privacy controls. Classify failures by responsible layer; do not conflate infrastructure retries with code passing on first attempt.

Verification: collect the configuration or runtime evidence supporting this control, identify an accountable owner, and record the exception path where the condition cannot presently be met.

\section{Adoption checklist}
\begin{itemize}
\item Identify applicable repositories, application owners and trust boundaries.
\item Define minimum pass criteria and measurable acceptance thresholds.
\item Review workflow permissions, external actions, outputs and retention.
\item Record exceptions with justification, owner and expiry.
\item Revalidate after action, runner or policy changes.
\end{itemize}
\end{document}
